\section{Conclusion}
\label{sec:conclusion}

This paper provides causal evidence that uninsured depositors discipline
community banks in response to a mandatory credit-risk disclosure.
Exploiting the CECL Day-One retained-earnings adjustment as an
information shock whose cross-sectional intensity was fixed by past
underwriting, we show that high-CECL community banks experienced
contractions in uninsured time deposits, simultaneous increases in
insured deposits, higher funding costs, lower profitability, and slower
loan growth relative to otherwise similar low-CECL peers.  The anchor
result is a within-bank composition shift: a triple difference with
bank $\times$ quarter fixed effects shows uninsured deposits falling
relative to insured deposits within the same bank and quarter in
proportion to the disclosed adjustment, an estimate that survives
allowing every balance-sheet characteristic implicated in the 2023
banking stress its own time path.  The displaced funding is traceable:
high-CECL banks replaced uninsured balances with insured brokered and
reciprocal deposits at higher cost, and part of the withdrawn money
reappears at local competitors.  The funding-cost and profitability
effects are amplified in markets with financially sophisticated
depositor bases.

The design addresses the three problems that have contaminated prior
tests of depositor discipline: simultaneity, accounting opacity, and
confounding from aggregate stress.  Together, our estimates are
difficult to attribute to anything other than informed uninsured
depositors acting on newly disclosed credit-risk information.

The policy implication is narrow and conditional.  Legislation that
would expand deposit insurance coverage for business accounts, such as
the Main Street Depositor Protection Act, would convert a share of the
uninsured funding we study into insured balances that, by design, carry
no incentive to monitor.  Our estimates identify the marginal cost of
that substitution on one channel only---the funding-cost, profitability,
and credit-supply consequences of a mandatory disclosure at community
banks in 2023---and do not measure the run-prevention benefits of
broader coverage or the general-equilibrium response of bank risk-taking.
A full policy evaluation requires a welfare framework outside the
present paper's scope.  What our evidence does support is a weaker
statement: proposals to expand coverage should be assessed for the
monitoring and pricing discipline they displace as well as for the
stability gains they deliver.
